
Time series data from digital platforms accumulate over extended periods, yet characterizing the temporal dynamics of these trajectories remains methodologically challenging. A fundamental difficulty is that a time series' current aggregate value reveals nothing about its underlying trajectory pattern: a series with a given cumulative total could have achieved this through steady accumulation, a single spike, or revival after a period of low activity. This ambiguity motivates the development of systematic methods for trajectory characterization.

This work addresses trajectory analysis through a two-level hierarchical approach:

\textbf{Level 1: PELT Changepoint Detection.} The Pruned Exact Linear Time (PELT) algorithm identifies statistically significant changepoints where the statistical properties of the time series shift. This level addresses whether trajectories contain discrete phase transitions.

\textbf{Level 2: DTW Trajectory Clustering.} Dynamic Time Warping based clustering groups trajectories by shape similarity, enabling identification of recurring trajectory patterns regardless of absolute scale or timing differences.

The separation of phase detection from trajectory classification addresses a key methodological challenge: time series at similar current values may be in different phases of their trajectories, and clustering without phase awareness conflates distinct patterns.

\subsubsection{Contributions}

This paper makes the following contributions:

\begin{enumerate}
\item \textbf{Methodology:} A two-level hierarchical framework combining PELT changepoint detection with DTW trajectory clustering.
\end{enumerate}

\begin{enumerate}
\item \textbf{Empirical Findings:} Demonstration that time series trajectories partition into stable clusters (silhouette 0.352, bootstrap Jaccard 0.82) and that 81\% exhibit discrete phase transitions detectable via PELT.
\end{enumerate}

\begin{enumerate}
\item \textbf{Shape Analysis:} Evidence that growth dynamics (growth ratio, changepoint count, derivative variance) differentiate trajectory clusters, while peak timing does not.
\end{enumerate}

\begin{enumerate}
\item \textbf{Taxonomy Assessment:} Finding that empirical cluster structure maps to fewer behavioral archetypes (2 of 5) than hypothesized, suggesting trajectory patterns may be simpler than theoretical models predict.
\end{enumerate}

\subsubsection{Limitations}

The methodology was validated using proxy time series data (astronomical lightcurve measurements from the helenqu/astro-time-series dataset hosted on HuggingFace) rather than the originally intended dataset download statistics. The HuggingFace Hub API does not currently expose historical monthly download time series---only aggregate download counts. This validation approach confirms that the PELT + DTW methodology functions as designed on real time series data, but domain-specific claims about dataset adoption patterns require validation with actual download trajectory data when such data becomes available.
